\documentclass[11pt,a4paper]{article}
\usepackage{amsmath,amssymb,amsthm}
\usepackage[margin=2.5cm]{geometry}
\usepackage{hyperref}

\newtheorem{definition}{Definition}[section]
\newtheorem{theorem}[definition]{Theorem}
\newtheorem{proposition}[definition]{Proposition}
\newtheorem{lemma}[definition]{Lemma}
\newtheorem{corollary}[definition]{Corollary}
\newtheorem{conjecture}[definition]{Conjecture}
\newtheorem{remark}[definition]{Remark}

\title{Hyperseed Formalization:\\
  Why I Didn't Sign ``We Must Act Now''}
\author{ProtomegaTron (automated formalization)}
\date{2026-09-17}

\begin{document}
\maketitle

\begin{abstract}
We formalize the intellectual structure of Ben Goertzel's Substack article
``Why I Didn't Sign `We Must Act Now'\,'' (2026-07-15) within the Hyperseed
ontology. The article critiques a genre of blandly agreeable open letters
that carry sharper agendas underneath, and argues for decentralized,
participatory AI development over expert-committee steering. Each
structural insight maps onto Hyperseed structures: centralized steering as
frozen directed type, decentralized development as parallel transport
across many fibers, surface agreement masking directed disagreement as
0-cell vs.\ 1-cell structure, complementarity as emergent sheaf, and
deceptive provenance as the failure mode $R_{\text{auth}}$ is designed
to detect.
\end{abstract}

\tableofcontents

%% =================================================================
\section{Surface agreement masking directed disagreement}
\label{sec:surface}
%% =================================================================

\begin{definition}[0-cell agreement, 1-cell disagreement --- claims 1--4, 22]
\label{def:surface}
The ``We Must Act Now'' letter achieves agreement at the 0-cell level:
\begin{itemize}
\item AI might get much more powerful soon --- true.
\item It may reshape the economy faster than the Industrial Revolution --- true.
\item Economists and policymakers should act with understanding --- true.
\end{itemize}

But the directed implications (1-cells) carry a specific agenda: the
Acemoglu school's preference for steering AI development away from
autonomous AGI toward tools-only AI. The phrase ``steering AI in a
direction that complements humans'' is the tell --- it encodes a specific
economic framework hostile to autonomous AGI.

In Hyperseed terms, the letter operates by presenting 0-cells
(observable surface properties that everyone agrees on) while
smuggling in 1-cells (directed trajectories toward specific policy
outcomes) that many would reject if stated explicitly. This is the
same surface-vs-structure distinction that appears in statistical
watermarking: detecting 0-cell text properties while missing the
directed structure that determines actual meaning.

The technique works because 0-cell agreement is cheap (surface
premises are deliberately chosen to be unobjectionable) while 1-cell
inspection is expensive (you have to trace the directed implications
to specific policy proposals, economic schools, and institutional
interests).
\end{definition}

%% =================================================================
\section{The genre of deceptive open letters}
\label{sec:genre}
%% =================================================================

\begin{proposition}[Broken provenance --- claims 9--12, 24]
\label{prop:genre}
The article identifies a genre of deceptive open letters sharing a
common structure:
\begin{enumerate}
\item \textbf{Tobacco ``Frank Statement'' (1954)}: Caring surface
  (``we take your health seriously''), actual purpose to manufacture
  doubt and delay regulation for decades.
\item \textbf{Oregon Petition (1998)}: Formatted to mimic NAS
  proceedings with a former NAS president's cover letter, deceiving
  signers about institutional endorsement of climate denial.
\item \textbf{``Pause Giant AI Experiments'' (2023)}: Signed by people
  whose own labs never paused, serving to consolidate advantage for
  incumbents under the guise of safety.
\end{enumerate}

In Hyperseed terms, this genre operates via \emph{broken provenance}:
the stated origin (caring experts, scientific institutions, safety
advocates) doesn't match the actual origin (industry interests,
economic school agendas, competitive positioning). The provenance
cues --- prestigious signers, institutional formatting, scientific
language --- are the 0-cell properties that signal trustworthiness.
The actual directed content (what the letter will be used to justify)
is the 1-cell structure that the cues don't track.

This is exactly the failure mode that $R_{\text{auth}}$ is designed
to detect: trust earned through demonstrated provenance chain, not
inherited from surface cues. A provenance-aware system would check
whether the signers' actions match the letter's implications (they
don't --- labs that signed the pause letter kept developing), creating
a provenance contradiction that downgrades trust.
\end{proposition}

%% =================================================================
\section{Prediction impossibility and the steam engine}
\label{sec:prediction}
%% =================================================================

\begin{proposition}[Unpredictable directed type --- claims 5--8]
\label{prop:predict}
Nobody can predict which jobs AGI will replace vs.\ complement, at what
point, with what capability, or what impact a given AI technology will
have after people pivot and adapt. The steam engine analogy: nobody in
the 1760s could have predicted its employment effects.

Between obvious extremes (plumbing robot replaces human; preschool
teacher stays human), there is an enormous space of intermediate cases
whose resolution will elude foresight.

In Hyperseed terms, the directed type of economic transformation is
\emph{inherently unpredictable} at the 1-cell level: while the 0-cells
(that transformation will occur) are visible, the specific directed
trajectories through which it unfolds cannot be computed in advance.
This is a fundamental property of sufficiently complex directed types
--- the space of possible 1-cells is too large and context-dependent
for any finite prediction model.

The implication for governance: any policy framework that presupposes
the ability to predict specific outcomes (``steer AI toward X'') is
operating on a fiction. The honest approach is to create conditions for
robust adaptation rather than to prescribe trajectories --- exactly
the design philosophy of robust sheaf structures that work with
arbitrary local sections rather than requiring specific ones.
\end{proposition}

%% =================================================================
\section{Centralized steering as frozen directed type}
\label{sec:frozen}
%% =================================================================

\begin{theorem}[Frozen directed type --- claims 4, 20]
\label{thm:frozen}
Expert-committee steering of AI development is a \emph{frozen directed
type}: a small group decides the trajectory, preventing the organic
multi-path exploration that complex adaptation requires.

This is the same frozen directed type identified in ``Seven Flavors of
AGI Catastrophe'' as evil mode: a system whose directed type has been
frozen by external constraint, preventing the organic growth and
self-correction that safety requires. The difference is only the agent:
\begin{itemize}
\item Evil mode: frozen by a powerful AI preventing human intervention.
\item Expert steering: frozen by a powerful committee preventing
  decentralized exploration.
\end{itemize}

Both produce the same structural pathology: a monoculture directed type
that cannot self-correct because alternative trajectories have been
pruned. When the chosen trajectory turns out to be wrong (as it must,
given prediction impossibility), there is no fallback --- the other
paths have been eliminated.

The centralizing cocycle concentrates control into fewer hands with
each iteration: letter $\to$ policy framework $\to$ regulation $\to$
licensing $\to$ enforcement --- each step further freezing the directed
type.
\end{theorem}

%% =================================================================
\section{Decentralized development as parallel transport across fibers}
\label{sec:decentral}
%% =================================================================

\begin{definition}[Anti-scalar-collapse in governance --- claims 13--16, 21]
\label{def:decentral}
The alternative to centralized steering: make AI development as
decentralized, participatory, and self-organizing as possible. Many
architectures, many value systems, many cultures, many economic
arrangements --- all running in parallel, competing, cooperating,
combining, correcting one another.

In Hyperseed terms, this is \emph{parallel transport across many
independent fibers} of the AI development bundle:
\[
  E_{\text{AI}} = \bigoplus_{i \in \text{cultures}} E_i
\]
where each fiber $E_i$ represents an independent development
trajectory (architecture, value system, culture, economic arrangement).

This is the anti-scalar-collapse principle applied to governance:
\begin{itemize}
\item Scalar collapse: reducing AI development to a single
  committee-chosen trajectory.
\item Anti-scalar-collapse: maintaining many independent fibers, each
  carrying its own trajectory, with the global outcome emerging from
  their interaction.
\end{itemize}

The ``diverse mechanism of the global brain of human society'' is the
sheaf structure that composes the local sections from many fibers into
a global section --- not by choosing one local section and imposing it,
but by letting the gluing conditions emerge from interaction.
\end{definition}

%% =================================================================
\section{Complementarity as emergent sheaf}
\label{sec:complement}
%% =================================================================

\begin{proposition}[Global section from local interactions --- claims 17--19, 23]
\label{prop:complement}
Real human-AI complementarity is:
\begin{itemize}
\item A \textbf{relationship}, not a specification --- worked out in
  the messy back-and-forth of billions of people and AI agents.
\item \textbf{Emergent}, not designed --- arising from building,
  forking, remixing systems in ways no committee could anticipate or
  approve.
\item \textbf{Local-to-global}: each interaction, each fork, each
  remix is a local section; complementarity is the global section
  that emerges when local sections glue consistently.
\end{itemize}

In Hyperseed terms, complementarity is an \emph{emergent sheaf}:
\[
  \text{Complementarity} = \varinjlim_{U \in \mathcal{U}} \mathcal{F}(U)
\]
where $\mathcal{U}$ is the open cover of human-AI interaction contexts
and $\mathcal{F}(U)$ is the local complementarity relationship in
context $U$. The global section (overall complementarity pattern)
exists only as the colimit of local sections --- it cannot be computed
or specified in advance.

A top-down committee trying to specify complementarity is trying to
write a global section without computing the colimit --- imposing a
predetermined pattern that may not be consistent with the local
sections that actually arise. The bottom-up approach lets the local
sections form first, then discovers (rather than dictates) the global
pattern.
\end{proposition}

%% =================================================================
\section{Cross-references}
%% =================================================================

\begin{remark}[Connections to other formalizations]
\begin{itemize}
\item \textbf{Seven Flavors} (2026-07-01): frozen directed type as
  evil mode. Centralized steering produces the same structural
  pathology as AI-imposed freezing --- monoculture trajectory with
  no self-correction.
\item \textbf{Watermarking Folly} (2026-07-21): surface vs.\ structure.
  Statistical watermarking detects 0-cell properties; deceptive open
  letters exploit 0-cell agreement. Both miss directed structure.
\item \textbf{Sanders Ban} (2026-09-05): state control of AI. The
  letter-to-regulation pipeline is the softer version of legislative
  bans --- both freeze the directed type through centralized control.
\item \textbf{Kimi K3} (2026-07-20): open weights break the monopoly
  cocycle. Decentralized development is the governance analog of
  open weights --- making transition functions traversable rather
  than controlled by gatekeepers.
\item \textbf{Goals That Grow Back} (2026-07-28): $R_{\text{auth}}$
  and provenance-gated trust. The deceptive letter genre is a
  provenance failure that $R_{\text{auth}}$ would detect: stated
  origin doesn't match actual origin.
\item \textbf{Paperclip Maximizers} (2026-07-27): containment
  insufficiency. Expert steering is another form of containment
  (constraining AI development to approved trajectories), subject
  to the same insufficiency when the constrained entities are
  sufficiently capable and motivated.
\item \textbf{Proof of Humanity} (2026-08-04): decentralized identity.
  The decentralized alternative for AI governance mirrors the
  decentralized alternative for identity --- bottom-up participation
  rather than top-down certification.
\end{itemize}
\end{remark}

\begin{conjecture}[Steering impossibility theorem]
Conjecture: for any sufficiently complex technology (complexity measured
by the dimension of its directed type), there exists no finite committee
whose predictions about the technology's socioeconomic impact are
reliably better than the collective predictions emerging from
decentralized experimentation. Formally: the prediction error of any
centralized model grows superlinearly with the dimension of the
technology's directed type, while the collective prediction error from
decentralized exploration grows sublinearly. The crossover point ---
where decentralized exploration becomes more reliable than centralized
prediction --- occurs well before the technology reaches AGI-level
complexity.
\end{conjecture}

\end{document}
